% =============================================================================
% Paper 2 Staging File: "A Geometric Observatory for Crisis Taxonomy"
% Content cut from Paper 1 during restructuring (2026-02-28)
% =============================================================================

% =============================================================================
% OBSERVATORY SUBSECTIONS (from Section 3: Geometric Observables)
% =============================================================================

\subsection{Fubini--Study Velocity}\label{subsec:velocity}

The Fubini--Study geodesic distance between consecutive quantum states:
\begin{equation}\label{eq:fs_velocity}
  v(t) = d_{\mathrm{FS}}\bigl(\psi(x_t), \psi(x_{t-1})\bigr)
       = \arccos\bigl(|\!\braket{\psi(x_t)}{\psi(x_{t-1})}\!|\bigr)\,.
\end{equation}
This measures the instantaneous rate at which the quasi-coherent state
traverses the projective Hilbert space $\CP^{n-1}$.  Rapid regime
transitions cause large state velocities.

\paragraph{Hyperparameters.}  Same as Berry.

\subsection{Speed Limit Ratio}\label{subsec:speed_limit}

The Mandelstam--Tamm bound constrains the maximum evolution speed of a
quantum state by the spectral gap $\Delta(t) = E_1(t) - E_0(t)$:
\begin{equation}\label{eq:speed_limit}
  r(t) = \frac{v(t)}{\Delta(t)}\,,
\end{equation}
where $v(t)$ is the FS velocity~\eqref{eq:fs_velocity}.  When
$r \to 1$, the system evolves at its maximum possible speed---a hallmark
of diabatic (sudden) transitions.  During financial crises, the
embedding dynamics approach the quantum speed limit, signaling that
the data manifold undergoes maximally rapid structural change.

\paragraph{Hyperparameters.}  Same as Berry.

\subsection{Dimensionality Collapse}\label{subsec:dim_collapse}

The inverse participation ratio (IPR) of the quantum metric eigenvalues
measures effective dimensionality collapse:
\begin{equation}\label{eq:ipr}
  \mathrm{IPR}(t) = \frac{\lambda_{\max}\bigl(g(x_t)\bigr)}
  {\mathrm{tr}\,g(x_t)}\,,
\end{equation}
where $g(x_t)$ is the quantum metric tensor~\eqref{eq:metric}.  When
$\mathrm{IPR} \to 1$, one eigenvalue dominates---the effective
dimensionality of the data manifold collapses to one.  This captures
herding behavior during crises, where all assets move in a single
correlated direction.

\paragraph{Hyperparameters.}  Same as Berry, plus optional subsampling
factor (default 5).

\subsection{Sectional Curvature Sign}\label{subsec:sectional_sign}

The fraction of negative sectional curvature in a rolling window:
\begin{equation}\label{eq:neg_frac}
  f_{\mathrm{neg}}(t) = \frac{1}{w} \sum_{s=t-w+1}^{t}
  \mathbf{1}\bigl[K(e_0, e_1; x_s) < 0\bigr]\,,
\end{equation}
where $K(e_0, e_1; x)$ is the sectional curvature in the principal
plane.  Negative sectional curvature indicates that nearby geodesics
diverge---analogous to chaotic dynamics.  A sustained high fraction
of negative curvature signals structural instability.

\paragraph{Hyperparameters.}  Same as Berry, plus window size $w$
for the rolling fraction (default 20) and optional subsampling.


% =============================================================================
% INTERACTION TEST (from Section 5: Results)
% =============================================================================

\subsection{Interaction Test}\label{subsec:interaction}

We classify crises \emph{a priori} as ``novel'' or ``conventional''
(Table~\ref{tab:crises}) and methods as ``geometric'' or
``classical.''  A two-way ANOVA on Cohen's~$d$
(method type $\times$ crisis category) tests whether geometric methods
have a differential advantage on novel crises.

We find no evidence of differential advantage:
$F_{\mathrm{interaction}}(1,116) = 0.38$, $p = 0.54$, partial
$\eta^2 = 0.003$.  Neither method type ($F = 3.33$, $p = 0.07$)
nor crisis category ($F = 3.85$, $p = 0.05$) reaches significance
at $\alpha = 0.05$.

The temporal stability documented in
Section~\ref{subsec:temporal_oos}---where Berry and QFI maintain
similar performance on post-2020 crises as pre-2020---is consistent
with \emph{additive} complementarity: geometric and classical methods
extract partially independent information, valuable for combination
regardless of crisis type.  This is precisely the condition under
which forecast combination is most
valuable~\cite{bates_granger1969}: independently informative
methods benefit from equal-weight averaging without requiring
crisis-type-specific allocation.


% =============================================================================
% LEARNED OPERATORS (from Section 5: Results)
% =============================================================================

\subsection{Learned Operators}\label{subsec:learned_operators}

The operator ablation above compares three \emph{heuristic}
construction methods.  We now ask whether \emph{learned} operators---fitted
by gradient descent to maximize crisis/non-crisis separability---can
improve upon the best heuristic choice.

\paragraph{Setup.}
Each operator $A_k$ ($8 \times 8$ Hermitian, $d^2 = 64$ real
parameters) is parameterized via diagonal and upper-triangle entries.
We optimize all $n_{\mathrm{op}} \times 64$ parameters jointly to
maximize Cohen's~$d$ between crisis and non-crisis raw scores, using
coordinate-wise numerical gradients (200 subsampled points per class,
150 steps, $\eta = 0.01$).  Two initializations are tested:
\emph{learned-from-random} and \emph{learned-from-PCA-inspired}.
Evaluation uses leave-one-crisis-out (LOCO) cross-validation: operators
are trained on all crises except one, then the held-out crisis is
scored through the \emph{full} detector pipeline (scaler $\to$ PCA $\to$
operators $\to$ rolling window $\to$ expanding z-score) to produce
effect sizes directly comparable to Table~\ref{tab:comparison}.

\begin{table}[htb]
\centering
\caption{Learned operator LOCO evaluation: median Cohen's~$d$ across
  12~crises for each operator method and detector.  Deltas ($\Delta$)
  show change from the corresponding heuristic baseline.}
\label{tab:learned_operators}
\small
\begin{tabular}{lcccc}
\toprule
Detector & Random & PCA-Insp. & Learned (Rand) & Learned (PCA) \\
\midrule
Berry Phase Rate    & \textbf{0.69}  & 0.40  & 0.64\,{\scriptsize($-0.05$)}  & 0.17\,{\scriptsize($-0.23$)}  \\
QFI Determinant     & 0.49  & 0.44  & \textbf{0.53}\,{\scriptsize($+0.04$)}  & 0.36\,{\scriptsize($-0.08$)}  \\
Multi-Lag Fidelity  & 0.38  & 0.38  & 0.38\,{\scriptsize($\pm0.00$)}  & \textbf{0.53}\,{\scriptsize($+0.15$)}  \\
\bottomrule
\end{tabular}

\medskip
\noindent{\footnotesize 150 gradient steps, $\eta = 0.01$, $h = 8$,
coordinate-wise numerical gradients with 10 subsampled parameters per step.
LOCO: train on 11 crises, evaluate on held-out through full detector pipeline.}
\end{table}

\paragraph{Results.}
Table~\ref{tab:learned_operators} reveals that gradient-learned operators
do not systematically improve over heuristic baselines.  For the Berry
Phase Rate detector, the best-performing method remains heuristic random
initialization (median $d = 0.69$); learning \emph{degrades} performance
in both cases ($-0.05$ and $-0.23$).  For QFI Determinant, learning
from random yields a modest improvement ($+0.04$), while for Multi-Lag
Fidelity, learning from PCA-inspired initialization produces the only
notable gain ($+0.15$).  Across all 36 detector--method combinations
(3~detectors $\times$ 12~crises), no learned method achieves a
statistically significant improvement over its baseline in a paired
comparison.

This negative result is informative: it suggests that the heuristic
operator constructions---particularly random initialization for Berry
Phase Rate---already capture the geometric structure relevant to regime
separation.  The high-dimensional optimization landscape ($8 \times
64 = 512$ parameters) may contain many local minima, and the
coordinate-wise gradient approach may be insufficient to find
substantially better solutions.  More expressive parameterizations
(e.g., neural operator networks) or alternative optimization strategies
(e.g., evolutionary methods) remain directions for future work.


% =============================================================================
% ONLINE DETECTION AND PnL BACKTEST (from Section 5 body + Appendix G)
% =============================================================================

\subsection{Online Detection and PnL Backtest}\label{subsec:online_backtest}

An exploratory online ensemble pipeline and PnL backtest are presented
in Appendix~\ref{sec:online_appendix}.  A na\"ive online ensemble
achieves AUC-ROC~$= 0.50$, but diagnostic analysis reveals that the
geometric observables exhibit \emph{lower} mean values during crisis
periods than normal periods (Mann-Whitney one-sided $p > 0.96$ for all
four features), explaining the near-chance performance.  The effect
sizes are meaningful (QFI: $d = 1.79$; spectral gap: $d = 1.48$)
but in the \emph{wrong direction} for standard anomaly detectors that
expect elevated scores during crises.

\paragraph{Signal inversion fix.}
Once the directional mismatch is diagnosed, the fix is straightforward:
negate features whose crisis-period mean is lower than the normal-period
mean.  With this correction, a simple expanding-percentile detector
achieves AUC-ROC~$= 0.84$---a 68\% improvement over the na\"ive
pipeline and a 36\% improvement over classical online CUSUM
(AUC-ROC~$= 0.62$).  An unsupervised variant (negating all features
without using crisis labels) achieves AUC-ROC~$= 0.84$ as well,
confirming that the directional correction is robust and does not
require labeled data.  A multi-scale fusion approach combining z-scores
at windows $w \in \{5, 10, 20, 40\}$ achieves AUC-ROC~$= 0.84$,
suggesting that the signal inversion is the dominant improvement
rather than multi-scale aggregation.

These results confirm that geometric observables \emph{do} carry
substantial online discriminative power, but the raw signal direction
differs from classical anomaly detectors: crisis periods correspond
to manifold \emph{contraction} (lower curvature, lower fidelity)
rather than expansion.  This finding is consistent with the
interpretation that crises represent collapse onto low-dimensional
attractors in the projective Hilbert space.

A long-flat strategy driven by the geometric signal matches ConstantVol
OOS Sharpe ($0.79$ vs.\ $0.79$) while reducing maximum drawdown by
$30\%$ (19.0\% vs.\ 27.0\%).  The walk-forward evaluation
(Section~\ref{subsec:walkforward}), which uses the offline z-score
detectors with expanding-window refits and monthly operator updates,
constitutes our primary causal deployment benchmark.


% =============================================================================
% APPENDIX G: ONLINE DETECTION (full appendix)
% =============================================================================

\section{Online Detection and PnL Backtest}\label{sec:online_appendix}

To evaluate practical deployability, we implement a strictly causal online
pipeline that computes geometric features and P(crisis) at each timestep
using only data available up to that point.  The pipeline comprises: (i)~an
\emph{OnlineGeometricFeatureComputer} that periodically refits
scaler/PCA/operators on a rolling 500-day window and outputs 10~features
(5~base geometric observables plus 5~EMA velocity features); (ii)~three
detectors---multivariate Bayesian filtering, periodic HMM refit, and
expanding-percentile RMS scoring---combined via a weighted ensemble
(weights $0.4/0.4/0.2$).

\begin{table}[htb]
\centering
\caption{Online detection performance (4,934 timesteps, 2005--2024).
  Detection rate and delay measured at 1~false alarm per year.}
\label{tab:online_detection}
\small
\begin{tabular}{lcccc}
\toprule
Detector & AUC-ROC & AUC-PR & Det@1/yr & Delay \\
\midrule
Online CUSUM                 & \textbf{0.619} & 0.330 & 67\% & 39d \\
Online Vol-Z                 & 0.576 & \textbf{0.366} & 67\% & 20d \\
Online Logistic (supervised) & 0.528 & 0.261 & 33\% & 82d \\
\textbf{Online Ensemble}     & 0.502 & 0.221 & 17\% & 20d \\
Online Percentile            & 0.502 & 0.218 &  0\% & --- \\
Online Bayesian              & 0.500 & 0.221 &  0\% & --- \\
\bottomrule
\end{tabular}
\end{table}

The na\"ive online results are sobering: the geometric ensemble achieves
only AUC-ROC~$= 0.502$, barely above chance.  Diagnostic analysis
reveals that geometric features have \emph{lower} mean values during
crises---the wrong direction for anomaly scoring.  The effect sizes are
meaningful (QFI: $d = 1.79$; spectral gap: $d = 1.48$) but inverted.

\paragraph{Signal inversion fix.}
Once the directional mismatch is identified, the fix is straightforward:
negate features whose crisis-period mean is lower than the normal-period
mean.  With this correction, the expanding-percentile detector achieves
AUC-ROC~$= 0.84$, a 68\% improvement over the na\"ive pipeline and 36\%
over classical CUSUM (0.62).  An unsupervised variant (negating all
features without labels) achieves the same AUC-ROC~$= 0.84$,
confirming the correction is robust.  The directional mismatch arises
because crises correspond to manifold \emph{contraction}---lower
curvature and fidelity---rather than the expansion assumed by standard
anomaly detectors.

\paragraph{PnL Backtest.}
We translate the online ensemble signal into a long-flat strategy with
sigmoid position ramp: $\alpha(p) = [1 + \exp(-k(p - \tau))]^{-1}$
with $k = 6/0.3 = 20$ and threshold $\tau = 0.5$.  The strategy targets
10\% annualized volatility, applies 0.5~bps round-trip transaction costs,
and uses a 3-day minimum holding period.

\begin{table}[htb]
\centering
\caption{Backtest results (2005--2024).  IS/OOS split at 2020-01-01.}
\label{tab:backtest}
\small
\begin{tabular}{lccccc}
\toprule
Strategy & Sharpe & IS Sharpe & OOS Sharpe & MaxDD \\
\midrule
ConstantVol SPY       & 0.63 & 0.57 & 0.79 & 27.0\% \\
Buy \& Hold SPY       & 0.51 & 0.45 & 0.68 & 56.5\% \\
\textbf{Geometric LongFlat} & 0.47 & 0.36 & \textbf{0.79} & \textbf{19.0\%} \\
Geometric LongShort   & 0.40 & 0.27 & 0.81 & 18.2\% \\
\bottomrule
\end{tabular}
\end{table}

The geometric strategy achieves comparable OOS Sharpe (0.79 vs.\ 0.79
for ConstantVol) while reducing maximum drawdown from 27.0\% to 19.0\%
($-30\%$ relative) and from 56.5\% to 19.0\% versus buy-and-hold
($-66\%$ relative).


% =============================================================================
% OBSERVATORY DISCUSSION CONTENT
% =============================================================================

% Observatory orthogonality paragraph (from Discussion > Interpretation):
%
% The four observatory extensions add further independent channels:
% Fubini--Study velocity captures instantaneous state movement, the
% speed limit ratio $v/\Delta$ detects diabatic transitions approaching
% the Mandelstam--Tamm bound, the dimensionality collapse index
% (metric IPR) measures herding-driven eigenspectrum concentration, and
% the sectional curvature sign fraction identifies sustained periods of
% geodesic divergence.  An adiabatic/diabatic taxonomy emerges: crises
% driven by sudden shocks (flash crashes, carry unwinds) produce high
% speed limit ratios, while slow structural crises (rate hikes, dot-com
% unwind) manifest as sustained dimensionality collapse.

% Observatory orthogonality paragraph:
%
% The seven geometric channels exhibit mean pairwise Spearman
% $|\rho| < 0.15$, substantially lower than the within-baseline
% correlation (typically $|\rho| > 0.3$).  This confirms that the
% geometric observatory provides genuinely independent detection
% information from a single embedding.  Negative results are equally
% informative: entanglement entropy ($d = 0.00$) and Ricci curvature
% rate ($d = 0.048$) provide negligible crisis separation despite
% theoretical plausibility, indicating that not all geometric
% quantities extracted from the QCML embedding are useful for regime
% detection.
